\subsection{Exponential and logarithm}

By the exponential power series we shall understand the element \(\sum_{n\geqslant 0}\frac{X^n}{n!}\) of \(\mathbf{Q}[[\mathbf{X}]]\) ; it will be denoted by \(\exp X\) or \(e^X\) .

PROPOSITION 13. --- In Q{[}{[}X, Y{]}{]} we have \(e^{X + Y} = e^X e^Y\) .

For the binomial formula gives

\[
\frac {(\mathbf {X} + \mathbf {Y}) ^ {n}}{n !} = \sum_ {i + j = n} \frac {\mathbf {X} ^ {i}}{i !} \frac {\mathbf {Y} ^ {j}}{j !}.
\]

Hence

\[
\begin{array}{r l} e ^ {X} e ^ {Y} & = \left(\sum_ {i \geqslant 0} \frac {X ^ {i}}{i !}\right) \left(\sum_ {j \geqslant 0} \frac {Y ^ {j}}{j !}\right) = \sum_ {i, j \geqslant 0} \frac {X ^ {i}}{i !} \frac {Y ^ {j}}{j !} = \sum_ {n \geqslant 0} \sum_ {i + j = n} \frac {X ^ {i}}{i !} \frac {Y ^ {j}}{j !} \\ & = \sum_ {n \geqslant 0} \frac {(X + Y) ^ {n}}{n !} = e ^ {X + Y}. \end{array}
\]

We shall define two elements \(e(\mathbf{X})\) , \(l(\mathbf{X})\) of \(\mathbf{Q}[[\mathbf{X}]]\) by

\[
e (\mathbf {X}) = e ^ {\mathbf {X}} - 1 = \sum_ {n \geqslant 1} \frac {\mathbf {X} ^ {n}}{n !}\tag{33}
\]

\[
l (\mathbf {X}) = \sum_ {n \geqslant 1} (- 1) ^ {n - 1} \frac {\mathbf {X} ^ {n}}{n}.\tag{34}
\]

We have

\begin{enumerate}
\def\labelenumi{(\arabic{enumi})}
\setcounter{enumi}{34}
\tightlist
\item
\end{enumerate}

\[
e (\mathbf {X} + \mathbf {Y}) = e (\mathbf {X}) + e (\mathbf {Y}) + e (\mathbf {X}) e (\mathbf {Y})
\]

\[
\mathrm{D} (e ^ {\mathrm{X}}) = \mathrm{D} (e (\mathrm{X})) = e ^ {\mathrm{X}}\tag{36}
\]

\[
\mathrm{D} (l (\mathbf {X})) = \sum_ {n \geqslant 0} (- \mathbf {X}) ^ {n} = (1 + \mathbf {X}) ^ {- 1}.\tag{37}
\]

PROPOSITION 14. --- We have \(l(e(\mathbf{X})) = e(l(\mathbf{X})) = \mathbf{X}\) .

The series l and e have no constant term and their terms of degree 1 are equal to X. By Prop. 10 of IV, p.~37 it suffices to prove the formula \(l(e(X)) = X\) . By the formulae (36) and (37) and Cor. 3 of IV, p.~33 we have

\[
\mathrm{D} (l (e (\mathrm{X}))) = (1 + e (\mathrm{X})) ^ {- 1} \mathrm{D} (e (\mathrm{X})) = (e ^ {\mathrm{X}}) ^ {- 1} e ^ {\mathrm{X}} = 1
\]

whence \(l(e(\mathbf{X})) = \mathbf{X}\) .

Let K be a Q-algebra, then the elements of K{[}{[}I{]}{]} without constant term form a commutative group E under addition. The elements of K{[}{[}I{]}{]} with constant term 1 form a commutative group M under multiplication (IV, p.~30). For each \(f \in E\) , we can define the elements \(e \circ f\) and \(l \circ f\) of E, and by Prop. 14 above, the mappings \(f \mapsto l \circ f\) and \(f \mapsto e \circ f\) are mutually inverse permutations of E; clearly they are continuous. Since \(\exp X = e(X) + 1\) , we see that the exponential mapping \(f \mapsto \exp f = e \circ f + 1\) is a continuous bijection of E onto M. By formula (4) of IV, p.~29 and Prop. 13, we have \(\exp(f + g) = (\exp f)(\exp g)\) for \(f, g \in E\) . Thus the exponential is an isomorphism of the topological group E onto the topological group M.

The inverse isomorphism of \(\mathcal{M}\) onto \(\mathcal{E}\) is called the logarithm and is written \(g \mapsto \log g\) . We thus have \(\log g = l(g - 1)\) for \(g\) in \(\mathcal{M}\) , and in particular,

\[
\log (1 + \mathbf {X}) = l (\mathbf {X}).\tag{38}
\]

Since the logarithm is a homomorphism of \(\mathcal{M}\) into \(\mathcal{E}\) , the formula \((1 + \mathbf{X})(1 + \mathbf{Y}) = 1 + (\mathbf{X} + \mathbf{Y} + \mathbf{X}\mathbf{Y})\) implies

\[
l (\mathbf {X}) + l (\mathbf {Y}) = l (\mathbf {X} + \mathbf {Y} + \mathbf {X Y}).\tag{39}
\]

Let \((u_{\lambda})_{\lambda \in L}\) be a summable family of elements of \(\mathcal{E}\) , then the family \((\exp u_{\lambda})_{\lambda \in L}\) is multipliable and we have

\[
\exp \left(\sum_ {\lambda \in L} u _ {\lambda}\right) = \prod_ {\lambda \in L} \exp u _ {\lambda}.\tag{40}
\]

Similarly, if \((f_{\lambda})_{\lambda \in L}\) is a multipliable family of elements of \(\mathcal{M}\) , the family \((\log f_{\lambda})_{\lambda \in L}\) is summable and we have

\[
\log \left(\prod_ {\lambda \in L} f _ {\lambda}\right) = \sum_ {\lambda \in L} \log f _ {\lambda}.\tag{41}
\]

Let \(g \in M\) , and let D be a continuous derivation of K{[}{[}I{]}{]}. We have \(\log g = l(g - 1)\) , hence by Cor. 3 of IV, p.~33 and (37) we have

\[
\mathrm{D} \log g = \mathrm{D} (g) / g.\tag{42}
\]

The expression \(D(g)/g\) is called the logarithmic derivative of g (relative to D).

